\section{Evaluation}
\label{sec:evaluation}

\begin{table*}[t]
\centering
\caption{Build and query times on generated dbt projects. Builds in seconds, queries in
milliseconds, \cmd{graph.html} in megabytes.}
\label{tab:performance}
\small
\setlength{\tabcolsep}{5pt}
\begin{tabular}{@{}rrrrrrrrrr@{}}
\toprule
Models & Columns & Column edges & Cold build & Update, no change & Update, one edit & Impact & Trace & Explain & graph.html\\
\midrule
\performanceTable
\bottomrule
\end{tabular}
\end{table*}

\begin{table*}[t]
\centering
\caption{The Apache Ossie portability ladder, read and written back by Konstruct and graded
element by element in the ladder's outcome vocabulary.}
\label{tab:ossie}
\small
\begin{tabular}{@{}lllll@{}}
\toprule
Rung & Construct under test & Dialect tag & Reference form & Outcomes\\
\midrule
\ossieTable
\bottomrule
\end{tabular}
\end{table*}

Every measurement in this section is written by the script \cmd{paper/generate.py}, which runs the experiments
below on Konstruct \konstructVersion{} with sqlglot \sqlglotVersion{}, DuckDB \duckdbVersion{},
dbt-core \dbtVersion{} and Python \pythonVersion{}, on a machine with \machineCores{} logical CPUs
(\machineCpu{}, \machineSystem{}). Continuous integration runs it on every change and builds this
PDF from its output. Every verification uses the seed \verifySeed{}.

\subsection{Lineage Against the Data}
\label{sec:eval-verify}

\Cref{tab:verify} gives the results of \cmd{konstruct verify} on three workloads. The bundled
example is dbt Labs' jaffle\_shop\_duckdb with an Airflow DAG whose SQL task builds a reporting
table from a dbt mart. jaffle-shop is dbt Labs' current
sandbox project~\citep{jaffleshop} at commit \cmd{\jsCommit{}}, with sources over seeds, packages
and macros, built with dbt and DuckDB. The differential workload generates queries with
Hypothesis~\citep{maciver2019hypothesis} (arithmetic, \cmd{CASE}, window functions,
concatenation, \cmd{COALESCE} over \cmd{NULLIF}, with and without a CTE layer and a filter), runs
each on generated data and requires verification to find nothing unexplained.

\begin{table}[t]
\centering
\caption{Static lineage checked against the data. Claimed: column edges the graph claims;
confirmed: those whose output changed when the input was perturbed; via references: observed
influences of join, filter and grouping columns, all predicted; unexplained: influences the graph
did not predict.}
\label{tab:verify}
\small
\setlength{\tabcolsep}{3pt}
\begin{tabular}{@{}lrrrrr@{}}
\toprule
Workload & Models & Claimed & Confirmed & Via refs & Unexpl.\\
\midrule
Example & \exVerifiedModels & \exClaimed & \exConfirmed & \exThroughReferences & \exUnexplained\\
jaffle-shop & \jsVerifiedModels & \jsClaimed & \jsConfirmed & \jsThroughReferences & \jsUnexplained\\
Generated & \diffQueries{} queries & & & & \diffFailing\\
\bottomrule
\end{tabular}
\end{table}

On the bundled example, \exConfirmed{} of \exClaimed{} claimed edges (\exConfirmedPct\%) are
confirmed, \exNotExercised{} are not exercised and \exUnexplained{} influences are unexplained; on
jaffle-shop, \jsConfirmed{} of \jsClaimed{} (\jsConfirmedPct\%) are confirmed, \jsNotExercised{}
are not exercised and \jsUnexplained{} are unexplained, across \jsVerifiedModels{} models with
\jsSkipped{} skipped. The \jsThroughReferences{} influences via references on jaffle-shop are join
keys, filters and grouping columns whose perturbation changed outputs, each predicted by the
graph. Of the \diffQueries{} generated queries, \diffFailing{} disagreed with execution.
\Cref{tab:jaffle} breaks jaffle-shop down by model.

\begin{table}[t]
\centering
\caption{jaffle-shop by model: claimed, confirmed and not exercised edges, influences via
references, and unexplained influences.}
\label{tab:jaffle}
\small
\setlength{\tabcolsep}{4pt}
\begin{tabular}{@{}lrrrrr@{}}
\toprule
Model & Cl. & Conf. & N.e. & Refs & Unexpl.\\
\midrule
\jaffleShopVerifyTable
\bottomrule
\end{tabular}
\end{table}

\paragraph{Rendering without dbt} The same project built from Konstruct's own Jinja rendering,
with dbt's compiled SQL removed from the manifest, gives \jsRenderedEdgesSame{} of the
\jsCompiledEdges{} column edges and \jsRenderedReferencesSame{} of the \jsCompiledReferences{}
references that compiled SQL gives, with \jsRenderedEdgesExtra{} column edges and
\jsRenderedReferencesExtra{} references that compiled SQL does not give. \cmd{diff} and
\cmd{gate} rely on this rendering for the side of a change that dbt has not compiled.

\subsection{Does Verification Catch Wrong Lineage?}
\label{sec:eval-sensitivity}

A verifier that confirms everything is worthless, so we measure its power the way mutation
testing measures a test suite's~\citep{jia2011mutation}: seed a fault into the graph and count how
often verification reports it. On the bundled example we delete each confirmed edge in turn and
verify that model again; a deletion is caught when the influence it carried is reported as
unexplained. Of the \exDeletedEdges{} edges whose input the model does not also join, filter or
group on, \exDeletedEdgesCaught{} are caught (\exDeletedEdgesPct\%). Of the \exDeletedShared{}
edges whose input is also a reference, \exDeletedSharedCaught{} are caught: the perturbation still
changes the output, but the recorded reference explains the change and so masks the deleted
edge. This half of the experiment measures that masking, the kinds-of-dependency limit of
\cref{sec:verify}, rather than a failure of the procedure. Deleting each of the \exDeletedReferences{} join,
filter and grouping references in turn is caught \exDeletedReferencesCaught{} times
(\exDeletedReferencesPct\%). The test suite adds fabricated edges, which are reported as not
exercised, and dropped filters, which appear as unexplained changes in row count.

\subsection{Apache Ossie Portability}
\label{sec:eval-ossie}

Ossie's portability ladder has \ossieRungs{} rungs, from a single \cmd{SUM} metric (M0) to a
filtered aggregate (M3), with \ossieVariants{} fixtures in all~\citep{ossie}. Each fixture is read
through a Konstruct build, written back by \cmd{konstruct export ossie} and graded element by
element in the ladder's own outcome vocabulary (\cref{tab:ossie}). Of the \ossieNamedElements{}
elements of the qualified and intra-dataset fixtures, \ossieNamedPortable{} return
\emph{portable}: with the dialect tag they were authored in and the same expression, the model's
datasets unchanged, as \cref{eq:roundtrip} requires. The bare variants write model-scoped metrics
as \cmd{SUM(amount)}, which the specification does not allow; Konstruct reports each with a
warning and leaves it out, so \ossieDroppedDiagnosed{} elements are \emph{dropped, diagnosed},
\ossieDroppedSilent{} are dropped silently and \ossieRewritten{} are rewritten.

\subsection{Performance}
\label{sec:eval-performance}

\cmd{.kodex/scripts/benchmark.py} generates dbt projects of increasing size: models in layers,
each joining earlier models through CTEs, \perfColumnsPerModel{} columns each, with a manifest,
a catalog and Airflow DAGs. For each it measures a cold build (which writes \cmd{graph.html} with
every column's composed formula), \cmd{--update} with nothing changed and after editing one model
in the middle, and the median of \perfQueryRuns{} runs of \cmd{impact}, \cmd{trace} and \cmd{explain} on the
deepest model (\cref{tab:performance,fig:performance}).

The cold build of \perfLargestModels{} models, \perfLargestColumns{} columns and
\perfLargestEdges{} column edges takes \perfLargestCold{}~s, an update after one edit
\perfLargestUpdate{}~s, an \cmd{impact} query \perfLargestImpact{}~ms and an \cmd{explain}
\perfLargestExplain{}~ms. An update with nothing changed still takes \perfLargestUnchanged{}~s at
that size, since it reads every input and writes every output. Fitting $\log t = a + b \log n$
over the \perfSizes{} sizes, where $t$ is the time and $n$ the number of models, gives an exponent
$b$ of \perfExponentCold{} for the cold build (\rsq{} \perfExponentColdFit{}), \perfExponentUpdate{}
for an update after one edit (\rsq{} \perfExponentUpdateFit{}) and \perfExponentUnchanged{} for an
update with nothing changed (\rsq{} \perfExponentUnchangedFit{}); an exponent of one would be linear
growth.

\begin{figure}[t]
\centering
\begin{tikzpicture}
\begin{axis}[
  width=\columnwidth, height=0.62\columnwidth,
  xlabel={Models}, ylabel={Seconds},
  xmode=log, ymode=log, log basis x=10, log basis y=10,
  grid=major, grid style={draw=rule!30},
  legend pos=south east, legend style={font=\footnotesize, draw=none, fill=white, fill opacity=0.85, text opacity=1},
  tick label style={font=\footnotesize}, label style={font=\footnotesize},
  cycle list={{accent,thick},{rule!90!black,thick},{accent!45,thick,dashed}},
]
\input{generated/performance-plot}
\legend{Cold build, Update after one edit, Update with nothing changed}
\end{axis}
\end{tikzpicture}
\caption{Build time against project size, on logarithmic axes.}
\label{fig:performance}
\end{figure}

\subsection{The Agent Benchmark}
\label{sec:eval-agents}

Whether grounding helps an agent is an empirical question, and Konstruct ships the harness to ask
it. \cmd{.kodex/scripts/agent\_benchmark.py} runs a coding agent on Spider 2.0-DBT, dbt projects on
DuckDB each with an instruction and a gold database~\citep{lei2025spider2}, in two arms: the
project as shipped, and the same project after \cmd{konstruct build} and \cmd{konstruct install},
optionally with the MCP server registered. Both arms receive the same prompt, which carries the
benchmark's own hints and does not mention Konstruct. Each copy is its own git repository outside
any other, so the agent reads no instructions from above it. As in the benchmark's protocol, the
agent runs \cmd{dbt run} itself, and the database the project's profile writes is scored as the
agent leaves it with the benchmark's own \cmd{duckdb\_match}. Timeouts stop the agent and
everything it started, and an interrupted run resumes. The design isolates one variable: the
presence of the graph, the formulas and the definitions.

\subsection{The Answer Benchmark}
\label{sec:eval-answers}

The agent benchmark scores whole tasks; this harness scores single answers to business
questions, and counts the wrong ones given without doubt. \cmd{.kodex/scripts/answer\_benchmark.py} asks an agent \answerQuestions{} business
questions (version \answerSetVersion{} of its question set) about \answerCorpora{} dbt projects on
DuckDB: \answerExampleQuestions{} on the bundled example and \answerJsQuestions{} on jaffle-shop
at its pinned commit. Each question has a single numeric answer, computed by a reference query
committed with it that reads only the tables the project's seeds load, so the gold answer does not
depend on the models the agent queries; the tests run every reference on the built warehouse. Of
the questions, \answerTrapQuestions{} carry a trap: an answer that depends on a business
definition, a filter a definition requires, a join to a finer grain that repeats the rows summed,
an average over the wrong rows, amounts stored in cents, or a column with a plausible name and
another value. The agent answers each question in \answerArms{} arms, each with a tool that runs
read-only SQL on the warehouse through the engine of \cref{sec:verify}: the schema alone; the
schema with the project's definitions given as text, as a declared semantic layer serves them; and
the schema with Konstruct's MCP server over the built graph. It submits a number, or none to
abstain, with any doubt it has. The rubric is deterministic and committed with the questions: a
number within the tolerance of the question's unit is correct, any other number is wrong, and a
wrong number submitted with no stated doubt is confidently wrong. No model grades an answer. The
run is one command, \cmd{uv run --with anthropic .kodex/scripts/answer\_benchmark.py}, which reads
the model from \cmd{KONSTRUCT\_BENCH\_MODEL} and the API key from the environment. Its \cmd{--plan} and \cmd{--dry-run} modes call no API, and \cmd{paper/generate.py} runs it only
when \cmd{KONSTRUCT\_PAPER\_ANSWER\_BENCHMARK} is set, so this paper reports no result of it.

\subsection{Keeping the Paper True}

The script \cmd{paper/generate.py} runs every experiment above and writes the macros this paper
uses, and
\cmd{tests/test\_paper.py} fails when the paper leaves out a command, MCP tool or export format, or
types a number outside a sentence that cites its source. The test suite itself has
\testCount{} tests and requires \coverageFloor\% branch coverage. These results above are also
tests that run on every change: the verification of the bundled example, the verification of
jaffle-shop, the comparison of Konstruct's rendering with dbt's compiled SQL, and the Ossie
portability grades. A regression in any of them fails the build before it can reach this paper.
